\documentclass[conference]{IEEEtran}
\IEEEoverridecommandlockouts
% IEEEtran sınıfı conference şablonudur; proje şartnamesinde IEEE formatı isteniyor.

\usepackage[utf8]{inputenc}
\usepackage[T1]{fontenc}
\usepackage[turkish]{babel}
\usepackage{cite}
\usepackage{amsmath,amssymb,amsfonts}
\usepackage{graphicx}
\usepackage{textcomp}
\usepackage{xcolor}
\usepackage{booktabs}
\usepackage{array}
\usepackage{hyperref}
\hypersetup{colorlinks=true, linkcolor=blue, citecolor=blue, urlcolor=blue}

\def\BibTeX{{\rm B\kern-.05em{\sc i\kern-.025em b}\kern-.08em
    T\kern-.1667em\lower.7ex\hbox{E}\kern-.125emX}}

\begin{document}

% ============================================================
\title{Derin Öğrenme Tabanlı Araba Gövde Tipi Sınıflandırma:\\
MobileNetV3 ile Transfer Öğrenme Yaklaşımı}

\author{
\IEEEauthorblockN{Ad Soyad\textsuperscript{*}}
\IEEEauthorblockA{
\textit{Bilgisayar Mühendisliği}\\
\textit{Kocaeli Üniversitesi}\\
Kocaeli, Türkiye\\
ogrenci.no@kocaeli.edu.tr}
\and
\IEEEauthorblockN{Ad Soyad}
\IEEEauthorblockA{
\textit{Bilgisayar Mühendisliği}\\
\textit{Kocaeli Üniversitesi}\\
Kocaeli, Türkiye\\
ogrenci2.no@kocaeli.edu.tr}
}

\maketitle

% ============================================================
\begin{abstract}
Bu çalışmada, sekiz farklı araba gövde tipini (SUV, VAN, Station Wagon, Micro,
Açık Tekerlekli, Sedan, Hatchback, Pick-Up) sınıflandıran derin öğrenme tabanlı
bir görüntü işleme sistemi sunulmaktadır. Sistem, ImageNet üzerinde önceden
eğitilmiş MobileNetV3-Large mimarisi üzerinde transfer öğrenme yaklaşımı ile
geliştirilmiştir. Veri seti çeşitli açık kaynaklardan toplanmış, sınıf dengesi
ve genelleme kapasitesini artırmak amacıyla veri çoğaltma teknikleri
uygulanmıştır. Model, 224$\times$224 girişle eğitilmiş; en iyi modelimiz
validation setinde \%XX doğruluk ve \%XX makro F1 skoru elde etmiştir.
Eğitim sürecinde EarlyStopping, Cosine Annealing öğrenme oranı planlayıcısı ve
sınıf dengesi için WeightedRandomSampler kullanılmıştır. Sistem, kullanıcıların
canlı tahmin yapabileceği Flask tabanlı bir web arayüzüne entegre edilmiştir.
\end{abstract}

\begin{IEEEkeywords}
araba sınıflandırma, derin öğrenme, transfer öğrenme, MobileNetV3,
PyTorch, konvolüsyonel sinir ağı, görüntü işleme
\end{IEEEkeywords}

% ============================================================
\section{Giriş}
Otomotiv endüstrisinde görüntü tabanlı araç sınıflandırma; trafik gözetim
sistemleri, otomatik park yönetimi, akıllı şehir uygulamaları ve ikinci el
araç pazarındaki içerik kategorilendirme gibi pek çok alanda kullanılmaktadır.
Bu görevin manuel yapılması maliyetli ve hatalara açıkken, derin öğrenme
yöntemleri yüksek doğrulukla otomasyon sağlamaktadır.

Bu çalışmada, görüntülerden sekiz farklı araba gövde tipinin doğru biçimde
sınıflandırılması hedeflenmiştir. Çalışmanın temel katkıları şunlardır:
\begin{itemize}
  \item Sekiz sınıf için açık kaynaklardan derlenmiş özelleştirilmiş bir
        veri setinin oluşturulması ve dengeli dağıtılması;
  \item Transfer öğrenme yaklaşımıyla MobileNetV3-Large mimarisinin
        problemin gereksinimlerine uyarlanması;
  \item Eğitim sürecinde overfitting riskini azaltan veri çoğaltma,
        label smoothing, weighted sampling ve early stopping tekniklerinin
        bütünleşik kullanımı;
  \item Kullanıcı dostu bir Flask web arayüzü ile canlı tahmin demosu.
\end{itemize}

% ============================================================
\section{İlgili Çalışmalar}
Görüntü sınıflandırma için ResNet \cite{resnet}, EfficientNet \cite{efficientnet},
ConvNeXt \cite{convnext} ve Vision Transformer (ViT) \cite{vit} gibi pek çok
derin mimari önerilmiştir. Mobil ve sınırlı kaynaklı cihazlarda çalışmak üzere
MobileNet ailesi geliştirilmiş; MobileNetV3 \cite{mobilenetv3} özellikle
NetAdapt arama algoritması ve Squeeze-and-Excitation \cite{senet} bloklarıyla
hız--doğruluk dengesinde öne çıkmaktadır. Bu çalışmada hem proje şartnamesindeki
95 MB model boyutu sınırı hem de arayüz tarafındaki tahmin hızı kriterleri
nedeniyle MobileNetV3-Large mimarisi tercih edilmiştir.

% ============================================================
\section{Materyal ve Yöntem}

\subsection{Veri Seti}
Veri seti birden fazla açık kaynaktan
(Kaggle\footnote{Cars Body Type Cropped ve Stanford Car Body Type Data
veri setleri referans alındı.}, Wikimedia Commons, ücretsiz stok foto siteleri)
manuel olarak derlenmiştir. Sınıf başına yaklaşık \textbf{XXX} görüntü
toplanmış olup dağılım Tablo~\ref{tab:dataset}'de verilmiştir.

\begin{table}[h]
\caption{Veri seti dağılımı}
\label{tab:dataset}
\centering
\begin{tabular}{lcc}
\toprule
\textbf{Sınıf}   & \textbf{Train} & \textbf{Validation} \\
\midrule
SUV              & XXX  & XX  \\
VAN              & XXX  & XX  \\
Station Wagon    & XXX  & XX  \\
Micro            & XXX  & XX  \\
Açık Tekerlekli  & XXX  & XX  \\
Sedan            & XXX  & XX  \\
Hatchback        & XXX  & XX  \\
Pick-Up          & XXX  & XX  \\
\midrule
\textbf{Toplam}  & \textbf{XXXX} & \textbf{XXX} \\
\bottomrule
\end{tabular}
\end{table}

Veri seti \%85 eğitim, \%15 doğrulama olarak rastgele bölünmüş; tekrarlanabilirlik
için sabit tohum (\texttt{seed=42}) kullanılmıştır. Test seti, proje gereği
sunum sırasında ders ekibi tarafından sağlanacak yeni görüntülerden oluşacaktır.

\subsection{Ön İşleme}
Tüm görüntüler $256\times256$'ya yeniden boyutlandırılmış, ardından merkezden
$224\times224$ kırpılmış ve ImageNet istatistikleriyle normalize edilmiştir:
\[\mu = (0.485,\,0.456,\,0.406),\quad \sigma = (0.229,\,0.224,\,0.225).\]

\subsection{Veri Çoğaltma (Data Augmentation)}
Eğitim sırasında genelleme yeteneğini artırmak için şu dönüşümler stokastik
biçimde uygulanmıştır:
\begin{itemize}
  \item \textbf{RandomResizedCrop}: $[0.75,1.0]$ ölçek aralığında rastgele kırpma;
  \item \textbf{RandomHorizontalFlip}: $p=0.5$ olasılıkla yatay yansıtma;
  \item \textbf{ColorJitter}: parlaklık, kontrast, doygunluk $\pm 0.2$, ton $\pm 0.05$;
  \item \textbf{RandomRotation}: $\pm 10^{\circ}$;
  \item \textbf{RandomErasing}: $p=0.25$ olasılıkla rastgele bölge silme
        (oklüzyona dayanıklılık).
\end{itemize}

\subsection{Model Mimarisi}
ImageNet üzerinde önceden eğitilmiş MobileNetV3-Large omurgası kullanılmıştır.
Orijinal sınıflandırıcı başı, 8 sınıflı problemimize uyarlamak için
yeniden tasarlanmıştır:
\begin{equation}
\mathrm{Linear}(960\!\to\!1280)\;\to\;\mathrm{Hardswish}\;\to\;
\mathrm{Dropout}(0.3)\;\to\;\mathrm{Linear}(1280\!\to\!8).
\end{equation}
Modelin toplam parametre sayısı yaklaşık 4.2 M, kaydedilen \texttt{.pt} dosya
boyutu $\approx$ 17 MB olup proje şartnamesinin 95 MB sınırının çok altındadır.

\subsection{Eğitim Konfigürasyonu}
Tablo~\ref{tab:hyper} eğitimde kullanılan tüm hiperparametreleri özetler.
\begin{table}[h]
\caption{Hiperparametreler}
\label{tab:hyper}
\centering
\begin{tabular}{ll}
\toprule
\textbf{Parametre} & \textbf{Değer} \\
\midrule
Optimizer            & AdamW \\
Backbone LR          & $1\times 10^{-4}$ \\
Classifier head LR   & $1\times 10^{-3}$ \\
Weight decay         & $1\times 10^{-4}$ \\
Scheduler            & Cosine Annealing \\
Batch size           & 32 \\
Epochs               & 40 (Early Stopping, patience=7) \\
Loss                 & CrossEntropy + Label Smoothing (0.05) \\
Class balancing      & WeightedRandomSampler \\
Mixed precision      & Etkin (AMP) \\
\bottomrule
\end{tabular}
\end{table}

Sınıf dengesizliğini azaltmak için WeightedRandomSampler ile her sınıftan
örneklenme olasılığı eşitlenmiş; label smoothing ile modelin aşırı güvenli
softmax çıkışları cezalandırılmıştır.

% ============================================================
\section{Deneysel Sonuçlar}

\subsection{Eğitim Süreci}
Şekil~\ref{fig:loss}'de eğitim ve doğrulama kayıp eğrileri,
Şekil~\ref{fig:acc}'de doğruluk eğrileri gösterilmektedir. Eğitim kaybı ile
doğrulama kaybı arasındaki farkın küçük tutulması, modelin aşırı uyum
göstermediğine işaret eder.

\begin{figure}[h]
\centering
\includegraphics[width=\linewidth]{plots/loss_curve.png}
\caption{Training \& Validation Loss eğrisi.}
\label{fig:loss}
\end{figure}

\begin{figure}[h]
\centering
\includegraphics[width=\linewidth]{plots/accuracy_curve.png}
\caption{Training \& Validation Accuracy eğrisi.}
\label{fig:acc}
\end{figure}

\subsection{Metrikler}
Tablo~\ref{tab:metrics}'de sınıf bazlı Precision, Recall ve F1-Score değerleri;
ortalama (macro ve weighted) skorlar ile birlikte sunulmuştur.

\begin{table}[h]
\caption{Sınıf bazlı performans metrikleri}
\label{tab:metrics}
\centering
\begin{tabular}{lccc}
\toprule
\textbf{Sınıf}   & \textbf{Precision} & \textbf{Recall} & \textbf{F1-Score} \\
\midrule
SUV              & 0.XX & 0.XX & 0.XX \\
VAN              & 0.XX & 0.XX & 0.XX \\
Station Wagon    & 0.XX & 0.XX & 0.XX \\
Micro            & 0.XX & 0.XX & 0.XX \\
Açık Tekerlekli  & 0.XX & 0.XX & 0.XX \\
Sedan            & 0.XX & 0.XX & 0.XX \\
Hatchback        & 0.XX & 0.XX & 0.XX \\
Pick-Up          & 0.XX & 0.XX & 0.XX \\
\midrule
\textbf{Macro avg}    & 0.XX & 0.XX & 0.XX \\
\textbf{Weighted avg} & 0.XX & 0.XX & 0.XX \\
\bottomrule
\end{tabular}
\end{table}

\subsection{Confusion Matrix}
Normalize edilmiş 8$\times$8 karışıklık matrisi
Şekil~\ref{fig:cm}'de gösterilmektedir. Diyagonal değerlerin yüksekliği
modelin doğru sınıflandırma başarısını yansıtır. Off-diagonal hücrelerdeki
yoğunlaşmalar, hangi sınıflar arasında karışıklık yaşandığını ortaya koyar
(örneğin Sedan--Hatchback geçişleri gibi görsel olarak benzer sınıflar).

\begin{figure}[h]
\centering
\includegraphics[width=\linewidth]{plots/confusion_matrix.png}
\caption{Normalize edilmiş Confusion Matrix.}
\label{fig:cm}
\end{figure}

\subsection{Çıkarım Süresi}
Modelin tek görüntü üzerindeki ortalama çıkarım süresi
\textbf{XX ms} (CPU) / \textbf{X ms} (GPU) ölçülmüştür; web arayüzünden
yapılan tahminlerde kullanıcı bekleme süresi 100 ms'nin altında kalmaktadır.

% ============================================================
\section{Web Arayüzü}
Eğitilmiş model, Flask tabanlı bir web uygulamasına entegre edilmiştir.
Arayüz şu bileşenleri içerir: (i) sürükle-bırak veya tıklayarak dosya yükleme,
(ii) yüklenen görselin canlı önizlemesi, (iii) tahmin edilen sınıfın güven
skoru ile büyük ve belirgin gösterimi, (iv) tüm sınıflar için olasılık
dağılımının yatay bar chart olarak animasyonlu sunumu ve
(v) çıkarım süresinin milisaniye cinsinden gösterimi. Tahmin sonuçları,
yüklenen orijinal görüntü ile yan yana düzende sergilenir.

% ============================================================
\section{Sonuç ve Gelecek Çalışmalar}
Bu çalışmada, transfer öğrenme yaklaşımıyla MobileNetV3-Large mimarisi
kullanılarak 8 araba gövde tipini sınıflandıran kompakt ($\approx$ 17 MB) ve
hızlı bir model geliştirilmiştir. Veri çoğaltma, label smoothing,
sınıf-dengeli örnekleme ve EarlyStopping teknikleri sayesinde modelin
genelleme kapasitesi artırılmıştır. Web arayüzü ile sistem son kullanıcıya
canlı tahmin imkânı sunmaktadır.

Gelecek çalışmalarda; veri seti boyutunun artırılması, daha güçlü
augmentation stratejileri (MixUp, CutMix), test zamanı augmentation (TTA)
ve Vision Transformer tabanlı modellerin karşılaştırılması planlanmaktadır.

% ============================================================
\begin{thebibliography}{00}
\bibitem{resnet}
K. He, X. Zhang, S. Ren, J. Sun,
``Deep Residual Learning for Image Recognition,''
\textit{CVPR}, 2016.

\bibitem{efficientnet}
M. Tan, Q. V. Le,
``EfficientNet: Rethinking Model Scaling for Convolutional Neural Networks,''
\textit{ICML}, 2019.

\bibitem{convnext}
Z. Liu et al.,
``A ConvNet for the 2020s,''
\textit{CVPR}, 2022.

\bibitem{vit}
A. Dosovitskiy et al.,
``An Image is Worth 16x16 Words: Transformers for Image Recognition at Scale,''
\textit{ICLR}, 2021.

\bibitem{mobilenetv3}
A. Howard et al.,
``Searching for MobileNetV3,''
\textit{ICCV}, 2019.

\bibitem{senet}
J. Hu, L. Shen, G. Sun,
``Squeeze-and-Excitation Networks,''
\textit{CVPR}, 2018.

\bibitem{pytorch}
A. Paszke et al.,
``PyTorch: An Imperative Style, High-Performance Deep Learning Library,''
\textit{NeurIPS}, 2019.

\bibitem{adamw}
I. Loshchilov, F. Hutter,
``Decoupled Weight Decay Regularization,''
\textit{ICLR}, 2019.
\end{thebibliography}

\end{document}
